\documentclass{practice}

\title{8}
\date{\today}

%\usepackage{etoolbox}
%\makeatletter
%\preto{\@verbatim}{\topsep=0pt \partopsep=-2\parsep }
%\makeatother

\usepackage{crysymb}

\usepackage{fancyvrb}
\fvset{listparameters=\setlength{\topsep}{0pt}\setlength{\partopsep}{0pt}}

\usepackage{listings}

\lstset{
basicstyle=\small\ttfamily,
columns=flexible,
breaklines=true
}

\begin{document}
\maketitle

\begin{task}{OpenSSL, why are you like this?}
  OpenSSL has two different tools for creating and verifying digital signatures:
  \begin{itemize}
    \item \href{https://docs.openssl.org/master/man1/openssl-pkeyutl/}{\texttt{pkeyutl}}
    \item \href{https://docs.openssl.org/master/man1/openssl-dgst/}{\texttt{dgst}}
  \end{itemize}

  I personally prefer the \texttt{pkeyutl} tool.
  The interface is a bit clearer, and semantically it makes more sense: signatures are an asymmetric primitive, and not a feature of hashing.
  However, many would disagree with me here, and the \texttt{dgst} tool does seem to be more popular in the wild, likely because it predates the existence of \texttt{pkeyutl}.

  For \texttt{pkeytul} take note of the \texttt{-rawin} and \texttt{-digest} options.
  Observe also the following notice for the \texttt{-digest} option:
  \begin{quotation}
    At this time, HashEdDSA (the ph or "prehash" variant of EdDSA) is not supported, so the -digest option cannot be used with EdDSA.
  \end{quotation}

  Read also the following info/warnings for \texttt{pkeyutl}:
  \begin{itemize}
    \item \href{https://docs.openssl.org/master/man1/openssl-pkeyutl/#ed25519-and-ed448-algorithms}{Note about Ed25519 and Ed448}
    \item \href{https://docs.openssl.org/master/man1/openssl-pkeyutl/#ec-algorithm}{Note about EC (applies to ECDSA)}
    \item \href{https://docs.openssl.org/master/man1/openssl-pkeyutl/#rsa-pss-algorithm}{Note about RSA-PSS}
  \end{itemize}

  For \texttt{dgst}, the \texttt{-sign} option has the following notice:
  \begin{quotation}
    Note that for algorithms that only support one-shot signing (such as Ed25519, ED448, ML-DSA-44, ML-DSA-65 andML-DSA-87) the digest must not be set. For these algorithms the input is buffered (and not digested) before signing. For these algorithms, if the input is larger than 16MB an error will occur.
  \end{quotation}
\end{task}

\begin{task}{Signing with \texttt{pkeyutl}}
  Generate a P-384 keypair for ECDSA, a keypair for Ed25519, and a keypair for RSA.

  \begin{Verbatim}
openssl genpkey -algorithm ed25519 \
  -out ed.priv.pem -outpubkey ed.pem
openssl genpkey -algorithm EC \
  -pkeyopt ec_paramgen_curve:P-384 \
  -out ec.priv.pem -outpubkey ec.pem
openssl genpkey -algorithm RSA \
  -pkeyopt rsa_keygen_bits:3072 -quiet \
  -out rsa.priv.pem -outpubkey rsa.pem
  \end{Verbatim}

  Create a message file
  \begin{Verbatim}
echo "Message" > message.txt
  \end{Verbatim}

  Sign each message twice using:
  \begin{itemize}
    \item Ed25519 in pure mode
    \begin{Verbatim}
openssl pkeyutl -sign \
  -in message.txt -inkey ed.priv.pem -out ed1.sig
openssl pkeyutl -sign \
  -in message.txt -inkey ed.priv.pem -out ed2.sig
    \end{Verbatim}
    and use \texttt{sha256sum ed*.sig} to verify that EdDSA is determinstic.

    \item ECDSA with the SHA-384 hash (since P-384 also provides 192 bits of security)
    \begin{Verbatim}
openssl pkeyutl -sign \
  -digest sha384 \
  -in message.txt -inkey ec.priv.pem -out ec1.sig
openssl pkeyutl -sign \
  -digest sha384 \
  -in message.txt -inkey ec.priv.pem -out ec2.sig
    \end{Verbatim}
    and use \texttt{sha256sum ec*.sig} to verify that ECDSA is probabilstic by default.

    Note that prior to OpenSSL v3.5, the \texttt{-rawin} parameter must be specified, if \texttt{-digest} is used.

    \item Deterministic ECDSA with the SHA-384 hash
    \begin{Verbatim}
openssl pkeyutl -sign \
  -digest sha384 -pkeyopt nonce-type:1 \
  -in message.txt -inkey ec.priv.pem -out detec1.sig
openssl pkeyutl -sign \
  -digest sha384 -pkeyopt nonce-type:1 \
  -in message.txt -inkey ec.priv.pem -out detec2.sig
    \end{Verbatim}
    and use \texttt{sha256sum detec*.sig} to verify that signing is deterministic.

    Deterministic ECDSA support was added to OpenSSL in \href{https://openssl-library.org/news/openssl-3.2-notes\#major-changes-between-openssl-31-and-openssl-320-23-nov-2023}{v3.2}.
    The parameter to enable it is obscure, and the documentation does not help you find it either\dots.

    \item RSASSA-PSS with the SHA-256 hash (3072-bit RSA keys provide 128 bits of security)
    \begin{Verbatim}
openssl pkeyutl -sign \
  -digest sha256 \
  -pkeyopt rsa_padding_mode:pss \
  -pkeyopt rsa_pss_saltlen:digest \
  -pkeyopt rsa_mgf1_md:sha256 \
  -in message.txt -inkey rsa.priv.pem -out pss1.sig
openssl pkeyutl -sign \
  -digest sha256 \
  -pkeyopt rsa_padding_mode:pss \
  -pkeyopt rsa_pss_saltlen:digest \
  -pkeyopt rsa_mgf1_md:sha256 \
  -in message.txt -inkey rsa.priv.pem -out pss2.sig
    \end{Verbatim}
    and use \texttt{sha256sum pss*.sig} to verify that PSS is probabilistic.

    \item RSASSA-PKCS1-v1\_5 with the SHA256 hash
    \begin{Verbatim}
openssl pkeyutl -sign \
  -digest sha256 \
  -pkeyopt rsa_padding_mode:pkcs1 \
  -in message.txt -inkey rsa.priv.pem -out pkcs1.sig
openssl pkeyutl -sign \
  -digest sha256 \
  -pkeyopt rsa_padding_mode:pkcs1 \
  -in message.txt -inkey rsa.priv.pem -out pkcs1.sig
    \end{Verbatim}

    and use \texttt{sha256sum pkcs1*.sig} to verify that PKCS\#1 v1.5 is deterministic.
  \end{itemize}
\end{task}

\begin{task}{Verifying signatures with \texttt{pkeyutl}}
  Using the keys and signatures generated in the previous task, verify them with \texttt{pkeyutl}.
  Note that the \texttt{-pubin} parameter is necessary if you are using a public key file.
  For private key files, the tool automatically uses its public part for verification.

  \begin{itemize}
    \item Ed25519
    \begin{Verbatim}
openssl pkeyutl -verify \
  -in message.txt -pubin -inkey ed.pem -sigfile ed1.sig
    \end{Verbatim}

    \item ECDSA, both probabilistic and deterministic
    \begin{Verbatim}
openssl pkeyutl -verify -digest sha384 \
  -in message.txt -pubin -inkey ec.pem -sigfile ec1.sig
    \end{Verbatim}

    Note that you must specify the digest, since OpenSSL defaults to SHA-256.

    \item RSASSA-PSS
    \begin{Verbatim}
openssl pkeyutl -verify -digest sha256 \
  -pkeyopt rsa_padding_mode:pss \
  -pkeyopt rsa_pss_saltlen:digest \
  -pkeyopt rsa_mgf1_md:sha256 \
  -in message.txt -pubin -inkey rsa.pem -sigfile pss1.sig
    \end{Verbatim}

    \item RSASSA-PKCS1-v1\_5
    \begin{Verbatim}
openssl pkeyutl -verify -digest sha256 \
  -pkeyopt rsa_padding_mode:pkcs1 \
  -in message.txt -pubin -inkey rsa.pem -sigfile pkcs1.sig
    \end{Verbatim}
  \end{itemize}
\end{task}
\end{document}
